% Section 5.1, from lectures/L05/appendix/a_the_bipolar_transistor.md.

\section{The bipolar transistor, and the switch}
\label{c5:app:a}

A practical treatment with as little semiconductor physics as possible. A transistor is an exponential, a current gain, and a saturation voltage, and those three facts carry the whole of Part~2.

\subsection{One equation, and one ratio}
\label{c5:sec:a1}
\[
  I_C = I_S\left(e^{V_{BE}/V_T} - 1\right), \qquad I_B = \frac{I_C}{\beta}
\]

$I_S$ is around $10^{-14}$ A, $V_T$ is 26 mV, and $\beta$, also written $h_{FE}$, is between 50 and 300.

\textbf{That is the diode equation of Chapter~\ref{c4:ch:feedback}~(Chapter~\ref{c4:ch:feedback})}, with the same $V_T$ and $I_S$, because the base-emitter junction \emph{is} a diode. Everything Chapter~\ref{c4:ch:feedback} established carries over:

\begin{itemize}
  \item \textbf{60 mV per decade of current.} A transistor at 1 mA and 0.65 V needs 0.71 V for 10 mA.
  \item \textbf{The incremental resistance is $V_T/I$}, 26 ohm at 1 mA.
\end{itemize}
\textbf{That second one is the single most important number in this book.} From Chapter~\ref{c7:ch:smallsignal} it is $r_e$, the intrinsic emitter resistance, and every gain, input resistance and output resistance in Part~2 is written in terms of it. It arrives here as a property of a diode.

\textbf{The emitter carries both currents:}

\[
  I_E = I_C + I_B = I_C\left(1 + \frac{1}{\beta}\right)
\]

which at $\beta = 50$ is 2 per cent more than $I_C$. That 2 per cent is why this book writes $I_E \approx I_C$.

\subsection{Beta, and why this book assumes 50}
\label{c5:sec:a2}
$h_{FE}$ is the least trustworthy number on any datasheet. A 2N3904 is specified at 100 to 300 at 10 mA, and the same device at 0.1 mA and at 100 mA can differ by a factor of three.

\textbf{This book assumes $h_{FE} = 50$ everywhere,} below the specified minimum of most small-signal parts, for one reason:

\textbf{A design that works at $h_{FE} = 50$ works at $h_{FE} = 300$. The reverse is not true.}

Hence the rule running through Part~2: \textbf{no design may depend on the value of beta.} Where beta appears in a result, the design is arranged so the result is insensitive to it. There is exactly one exception, the input resistance of an emitter follower in Chapter~\ref{c8:ch:follower}, and the chapter says so.

\subsection{Three regions}
\label{c5:sec:a3}
Two junctions, each forward or reverse biased, give four combinations, of which three are used.

\begin{booktable}{@{}lllL@{}}
  \thead{Region} & \thead{Base-emitter} & \thead{Base-collector} & \thead{Behaviour} \\
  \midrule
  Cutoff & Reverse & Reverse & No current. An open switch. \\
  Forward active & Forward & Reverse & $I_C = \beta I_B$. An amplifier. \\
  Saturation & Forward & Forward & $I_C < \beta I_B$. A closed switch. \\
\end{booktable}

The fourth, reverse active, is a transistor used backwards; it works badly and nothing uses it on purpose.

\bookfigure{lectures/L05/appendix/images/bjt_output.png}{Collector current against collector-emitter voltage for five base currents, with the fifty ohm load line drawn across them from five volts on the horizontal axis to a hundred milliamps on the vertical, and the two operating points of a switch marked at the ends of that line.}{c5:fig:bjtoutput}

The curves are what the transistor allows, the straight line what the external circuit allows, and the operating point is where they cross. \textbf{A switch uses only the two ends of that line;} everything between them is where an amplifier lives and where a switch dissipates for nothing.

\textbf{Saturation is where the model stops being simple.} In the active region $I_C$ does not depend on $V_{CE}$, so the transistor is a current source. Below about 0.2 V it does, steeply, so it is a resistor instead. The simple model draws that knee as a corner and a real device rounds it off.

\subsection{Designing a switch}
\label{c5:sec:a4}
A switch must be \textbf{on hard}, meaning saturated, meaning a base current above $I_C/\beta$. Since $\beta$ may not be trusted, it is chosen from a \textbf{forced beta} instead:

\[
  \beta_{forced} = \frac{I_C}{I_B}, \qquad \text{chosen as 10}
\]

Ten times the base current the active region would need guarantees saturation for any device with $h_{FE}$ above 10, which is every silicon transistor ever made.

\bookfigure{lectures/L05/appendix/images/switch.png}{A switch: five volt logic through a base resistor into an NPN transistor, whose collector drives a fifty ohm load to a five volt rail and whose emitter is grounded, with annotations giving the base resistor as four hundred and seventy ohms and noting that the saturated collector-emitter voltage is not zero.}{c5:fig:switch}

For 100 mA of load current from a 5 V logic level:

\[
  I_B = \frac{100\ \text{mA}}{10} = 10\ \text{mA}, \qquad
  R_B = \frac{5 - 0.7}{10\ \text{mA}} = 430\ \Omega
\]

so 470 ohm from the E12 series, giving 9.1 mA and a forced beta of 11.

\textbf{Note the 0.7 V, where every bias point in this book uses 0.65.} A forced beta of 10 drives the base ten times harder than the active region would, and 60 mV a decade is what that costs: the junction sits above the small-signal figure, and \code{baseResistor} uses 0.7 for exactly this reason. At larger currents it moves again, which is why Exercise~\ref{c5:ex:5} needs 0.8.

\textbf{What forced beta costs.} A tenth of the load current thrown away in the driver, and stored base charge that must be removed before the device turns off. A hard-saturated transistor is slow to turn off for exactly that reason, and a fast switch is deliberately kept out of deep saturation.

\textbf{What the saturation voltage costs.} 0.1 to 0.3 V across the device, so 10 to 30 mW at 100 mA. A linear stage delivering the same current at half the supply dissipates 250 mW. \textbf{That ratio is the entire reason switching exists.}

\subsection{What the model here is blind to}
\label{c5:sec:a5}
\begin{itemize}
  \item \textbf{Bulk resistance.} The ideal transport model gives a saturation voltage of about 57 mV for the switch above; a real device gives 100 to 300 mV. \textbf{The model is optimistic, and the direction matters:} a design relying on the modelled figure runs hotter than predicted.
  \item \textbf{Speed.} Charge storage, transit time and junction capacitance are all absent, so nothing here says how fast a transistor switches. That is usually what decides the choice of part.
  \item \textbf{Breakdown.} Above some collector-emitter voltage the device conducts whatever the base does, and an inductive load will drive it there unless something is done about that.
  \item \textbf{Second breakdown.} A power transistor can fail at combinations of voltage and current that are individually within specification. Nothing here models it.
\end{itemize}
